%% ---------------------------------------------------------------------
\section{DIS 的扭度 4：从矩到 \texorpdfstring{$x$}{x} 空间}\label{sec:power-twist4}

单举 DIS 是幂次修正唯一有完整算符结构的强子过程。对偶数 $n$，OPE 给出
\begin{equation}\label{eq:power-moments-tw4}
M_2^{(n)}(Q^2)\equiv\int_0^1\!\dd x\,x^{n-2}F_2
=\sum_iC^{(n)}_{2,i}A^{(n)}_i+\frac1{Q^2}\sum_jC^{(n),\tau=4}_{2,j}A^{(n),\tau=4}_j+O\Big(\frac{1}{Q^4}\Big),
\end{equation}
第二项来自扭度 4 的夸克—胶子与四夸克定域算符（\cref{prop:dis-twist4}，\tagR，$R_{\mathrm{OPE}}$）；扭度 2 算符的迹项给出另一类、纯运动学的靶质量修正（\cref{prop:dis-tmc}）。按\cref{lem:dis-twist-classification}的共线扭度计数，强子与硬部分之间每多一条横向部分子就多一个 $1/Q$（\cref{fig:higher-twist}）；扭度 3 项在非极化 DIS 中没有可用的横向矢量而为零，故第一个修正是 $1/Q^2$。

固定 $n$ 的扭度 4 算符是一族（$n-2$ 个 $D^+$ 可分配到不同的场上）。对所有 $n$ 重新求和，得到场位于同一条光锥直线 $z=(0^+,z^-,\mathbf0_T)$ 的不同点上、相邻场由直 Wilson 线（\cref{def:prelim-wilson}）相连的规范不变算符；其矩阵元关于各光锥分离的 Fourier 变换即扭度 4 关联函数 $T^{(4)}_i([x])$，$[x]$ 为独立的动量分数。典型代表是 Ellis--Furmanski--Petronzio “横向基”中的 $\bar\psi\gamma^+D_\perp^jD_{\perp j}\psi$ 型关联与四夸克关联（其余协变导数已用运动方程消去）\cite{EllisFurmanskiPetronzio1982,EllisFurmanskiPetronzio1983}；在 $A^+=0$ 规范下它们是只依赖于纵向动量分数的多部分子分布\cite{Jaffe1983}。

\begin{proposition}[DIS 扭度 4 修正的 $x$ 空间形式 \tagA]\label{prop:power-dis-twist4-x}
设 $0<x<1$ 固定，$W^2\gg\Lambda^2$。当 $Q\to\infty$ 时，强子靶的非极化结构函数满足
\begin{equation}\label{eq:power-twist4-x}
F_2(x,Q^2)=F_2^{\mathrm{tw2}}(x,Q^2)+\frac{1}{Q^2}\sum_i\int\!\dd x_1\,\dd x_2\,H_i\Big(x_1,x_2,x;\frac{Q^2}{\mu^2},\as(\mu)\Big)\,T^{(4)}_i(x_1,x_2;\mu)+R_4(x,Q),
\end{equation}
其中 $F_2^{\mathrm{tw2}}$ 是扭度 2 的贡献（\cref{thm:disall-factorization}的领头项，加上由扭度 2 矩阵元决定的靶质量修正，其 $x$ 空间形式与处方有关，\cref{prop:dis-tmc}），$T^{(4)}_i$ 是一组完备的扭度 4 光锥关联函数（对四场关联函数，$\int\dd x_1\dd x_2$ 理解为对全部独立动量分数的积分），$H_i$ 是与靶无关的微扰硬函数，$R_4=O(\Lambda^4/Q^4)$（至多差对数），系数依赖于 $x$，在 $x\to1$、$x\to0$ 时不均匀。并且：(a) $1/Q^2$ 项的矩 $\int_0^1\dd x\,x^{n-2}(\cdots)$ 等于\eqref{eq:power-moments-tw4}的扭度 4 项；(b) 对核靶，$T^{(4)}_i$ 含按 $A^{1/3}$ 增强的双重散射部分（\cref{prop:disall-nuclear}）。
非微扰地位：\cref{asm:intro-bridge}（MA1、MA2、MA4）。
% errtype: E-S
\uses{prop:dis-twist4,prop:dis-tmc,thm:disall-factorization,asm:intro-bridge}
\end{proposition}

\begin{proofsketch}
1.~由\cref{prop:dis-twist4}，每个偶数矩的 $1/Q^2$ 项是扭度 4 定域算符的贡献。2.~自旋 $n$ 的扭度 4 定域算符是光锥算符在各分离点处 Taylor 展开的系数，故 $A^{(n),\tau=4}_j$ 是 $T^{(4)}_i$ 的多重矩的线性组合；\eqref{eq:power-twist4-x}对 $x$ 取矩即得 (a)。反之，固定 $n$ 的矩只确定 $H_i$ 的有限多个矩，推不出 $x$ 空间公式。3.~（共线展开）在割图中把硬部分对部分子横动量展开到二阶，加上额外一两条共线胶子接入硬部分的图，用 Ward 恒等式组合成规范不变的横向协变导数，得到“硬函数 $\otimes$ 关联函数”的形式\cite{EllisFurmanskiPetronzio1982,EllisFurmanskiPetronzio1983}；树图硬函数见\cite{Jaffe:1982pm,EllisFurmanskiPetronzio1983}，使规范不变性显然的推导见\cite{Qiu:1988dn}。4.~（只被论证的一步）在 $1/Q^2$ 阶，主导区域恰是领头区域展开到二阶与额外共线部分子接入硬部分的区域，软区与 Glauber 区不给出新的贡献（单个强子，末态单举求和，参见\cref{thm:disall-factorization}）；此步只给出 E-S 型陈述。完整论证见\cite{EllisFurmanskiPetronzio1983,Qiu:1988dn}。
\end{proofsketch}

\paragraph{地位与缺口}
矩空间的结构是 $R_{\mathrm{OPE}}$ 严格的；$x$ 空间形式为 \tagA，最弱的一步是第 4 步的次领头幂次区域分析（E-S 型）——它相当于把\cref{asm:intro-bridge}(Br1) 推广到次领头幂次，而编号假设只覆盖领头幂次。$F_2^{\mathrm{tw2}}$ 经由\cref{thm:disall-factorization}（\tagA）继承 \cref{asm:appN-collins}、\cref{asm:appN-envelope}、\cref{asm:appN-separation}、\cref{asm:appN-cutsum}、\cref{asm:appN-ward-product}、\cref{asm:appN-uv}、\cref{asm:appN-multren}，以及该论证中被明确假定的一步（\cref{asm:appN-brst-nonlocal}，[C]；\cite[第~11.3.2~节]{Collins2011}）。此外，扭度 4 与领头扭度通过重整化子歧义混合（\cref{conj:power-renormalon-ht}），扭度 2 系数只在相对 $O(\Lambda^2/Q^2)$ 内有意义（\cref{prop:power-rearrangement}）。升级途径：次领头幂次区域分析的积分范数估计，以及\cref{thm:disall-factorization}本身的升级。

\begin{example}[扭度展开在 $x\to1$ 处的增强]\label{ex:power-twist4-largex}
设 $q(x)=(1-x)^b$，$M_n[q]\equiv\int_0^1x^{n-1}q=B(n,b+1)$。修正 $\delta q=c\,b\,(1-x)^{b-1}$ 的矩为 $c\,b\,B(n,b)=c\,(n+b)M_n[q]$，所以“矩的修正随 $n$ 线性增长”与“相对修正 $\delta q/q=c\,b/(1-x)$”是同一件事：取 $c\sim\Lambda^2/Q^2$，大 $x$ 处的实际展开参数是 $\Lambda^2/[Q^2(1-x)]$，即真正的硬标度是 $Q\sqrt{1-x}$（$W^2\simeq Q^2(1-x)/x$）。扭度 4 的“重整化子模型”正有这种行为：大 $\beta_0$ 近似中 $u=1$ 的重整化子留数给出的 $F_2$ 形状含 $-4/[1-x]_+$ 与 $-\delta'(1-x)$\cite[式~(5.33)]{Beneke1999}，其矩分别为 $4S_1(n-1)\sim4\ln n$ 与 $-(n-1)$；与单胶子线相联系的高阶重整化子修正按 $[\Lambda^2/(Q^2(1-x))]^n$ 增长\cite[式~(5.37)]{Beneke1999}（参见\cite{BenekeBraun1995b,DasguptaWebber1996}）。该模型只给形状不给归一化，而拟合所得的“扭度 4 修正”随扭度 2 项微扰阶数的提高而减小\cite[第~5.2.4~节]{Beneke1999}：二者不能被干净地分开。（Mellin 关系均用符号计算核对过。）
\uses{prop:power-dis-twist4-x}
\end{example}
